\section{Chiral Protocols}

\subsection{The Two Valid Routing Protocols: Singlets and Fractions}

Having established that Electric Charge is the geometric footprint of an information packet crossing the macroscopic boundary, we apply the Unitary Flux Constraint (information cannot be created or destroyed). The finite network permits exactly two ways to route a payload without triggering Causal Aliasing:

\paragraph{Singlet Routing (The Lepton Constraint):}
If an information packet bypasses the $3$D internal color bulk entirely, $100\%$ of its payload routes directly through the $\chi=2$ topological boundary. Because there is no internal geometric dilution, the topological footprint remains whole. The phase flux crossing the boundary is perfectly conserved as a single, undivided integer:
\begin{equation}
    Q_{lepton} = \frac{Y}{1} = -1
\end{equation}
\textbf{Conclusion:} Leptons are the non-diluted routing protocol. Bypassing the internal friction of the color bulk explains their macroscopic stability as point-like, integer-charged particles.

\paragraph{Fractional Routing (The Quark Constraint):}
If a packet requires the internal color bulk, it must route its data through the $N_{res} = 3$ internal gluon links. To preserve $SU(3)$ interaction symmetry, the flux must distribute isotropically across all three internal dimensions. When this colored packet simultaneously anchors to the spacetime manifold, the flux passing through the boundary is geometrically diluted by this 3-fold internal distribution:
\begin{equation}
    Q_{quark} = \frac{Y}{\sigma - \chi} = \frac{Y}{3}
\end{equation}
\textbf{Conclusion:} Quarks do not simply ``happen'' to have fractional charges. Fractional charge is the uniquely required, impedance-matched routing protocol for colored packets. If a state occupying the color channels possessed an integer charge, its flux would asymmetrically overload the destination node's boundary channels, violating Unitarity and triggering immediate decay.

\subsection{The Four Base Protocols}

Having established the strict constraints of transit (the Phase Offset, the Internal Embedding, and the Boundary Flux), we can apply the Unitary Flux Constraint (information cannot be created or destroyed) to exhaustively enumerate all mathematically valid ways to route a phase-tracked (Spin-$1/2$) payload across the $\nu=16$ active capacity channels. This results in the lattice permitting exactly \textbf{four distinct base routing protocols}. From these four protocols  the fractional quarks and integer leptons of the Standard Model appear.

A base protocol is defined by two geometric binary choices: \textit{Does it route through the internal Color Bulk?} and \textit{What is its topological boundary footprint?}

\begin{enumerate}
    \item \textbf{The Electron Protocol (Boundary-Coupled Singlet):} The payload bypasses the internal color bulk entirely, routing $100\%$ of its flux directly through the $\chi=2$ topological boundary. Because there is no internal geometric dilution, the topological footprint is a whole, undivided integer ($Q=-1$). It requires a base anchor of $1$.
    
    \item \textbf{The Up Protocol (Color-Coupled Singlet):} The payload routes through the $N_{res}=3$ internal color bulk. To maintain unitary phase matching at the destination node, the boundary flux must be geometrically diluted by this 3-fold internal distribution ($Q=+2/3$). Because it utilizes the internal bulk, its flux tubes must strictly anchor to the $D=4$ spacetime manifold, incurring a base anchor of $D=4$.
    
    \item \textbf{The Down Protocol (Color-Coupled Singlet):} A fractional routing protocol ($Q=-1/3$) that shares the $D=4$ manifold anchor of the Up protocol, but carries the additional static geometric load of the internal interaction symmetry ($+\sigma$). 
    
    \item \textbf{The Neutrino Protocol (Unanchored Singlet):} The payload bypasses the color bulk (Colorless) and possesses zero boundary flux footprint ($Q=0$). Lacking both the $D$ manifold anchor and the $\chi$ boundary anchor, it cannot form a localized standing wave (a Knot). It must propagate as a delocalized lattice vibration (a Phonon).
\end{enumerate}

\paragraph{The Geometric Origin of CPT Symmetry:}
In standard physics, the CPT theorem is an axiomatic requirement of Lorentz-invariant quantum field theories. In Informational Energetics, it emerges as a single, rigid geometric lock on the $E_8$ substrate:
\begin{itemize}
    \item \textbf{Charge (C):} The $\chi=2$ spherical boundary strictly admits two topological orientations (inside/outside).
    \item \textbf{Parity (P):} The double-cover $Spin(8)$ transition matrix rigorously isolates two chiralities (Left/Right).
    \item \textbf{Time (T):} The causal metric projection assigns active temporal updates exclusively to the Left-Handed sector.
\end{itemize}
Because the topological boundary ($\chi$), the phase rotation ($Spin$), and the temporal generator ($\Delta$) are structurally bound to the exact same lattice hardware, a reversal of one necessitates the reversal of all to maintain a valid, persistent routing protocol. CPT is inviolable because the substrate geometry is inviolable.

\subsection{Chiral Routing and the Neutrino Exception}

Having identified the four base routing protocols (Up, Down, Electron, Neutrino), we can map them to the active chiral state space ($\nu=16$). On the $E_8$ substrate, the metric projection assigns active temporal updates to the Left-Handed sector, while the Right-Handed sector remains statically embedded in the internal symmetry space (as derived in Paper I). For an information packet to manifest as both a Left-Handed active state and a Right-Handed static state (creating the chiral basis of a massive particle), the protocol must possess a \textbf{Geometric Anchor} capable of bridging these orthogonal spaces. 

\begin{itemize}
    \item \textbf{Quarks} possess the $D=4$ Manifold Anchor, required to route their flux through the 3D internal color bulk. This anchor bridges the chiral sectors.
    \item \textbf{Electrons} possess the $\chi=2$ Topological Anchor, required to define their continuous $Q=-1$ boundary flux footprint. This anchor bridges the chiral sectors.
\end{itemize}

Because they possess geometric anchors, the Up, Down, and Electron protocols can maintain both Left-Handed and Right-Handed routing states. 

\paragraph{The Neutrino Duality: Why it Lacks a Right-Handed State}

The Neutrino protocol is uniquely devoid of geometric anchors. It is colorless (bypassing the $D=4$ color bulk anchor) and electrically neutral (possessing zero boundary flux footprint on $\chi=2$). 

Because it lacks any geometric mechanism to anchor itself to the macroscopic spacetime manifold, the Neutrino \textbf{cannot form a localized standing wave}. Without an anchor to bridge the active and static sub-lattices, it cannot reflect into the Right-Handed internal sector. 

This provides a rigorous geometric proof for two distinct phenomena:
\begin{enumerate}
    \item \textbf{The Phonon Duality:} The Neutrino is a Fermion because it propagates strictly through the chiral spinor channels, but it acts as a delocalized Phonon (following the inverse resonance hierarchy $M \propto \Delta^{-n}$) because it possesses zero topological anchors to halt its propagation.
    \item \textbf{Chiral Asymmetry:} The Neutrino protocol exists strictly as a Left-Handed active routing state ($\nu_{eL}$). A Right-Handed neutrino routing protocol ($\nu_R$) is geometrically forbidden. 
\end{enumerate}

\paragraph{The Sterile Vacancy: The Shannon Capacity Lock}

If the Right-Handed neutrino is geometrically forbidden as a routing protocol, what occupies its expected channel space? 

As established in the \textit{Capacity Partitioning} rules (Paper III), the substrate operates under finite Shannon capacity limits. To maintain unitary phase synchronization when the active channels approach saturation, the network must lock exactly one channel as a sterile reference state (analogous to a static pilot tone in dense digital signaling). 

This requirement perfectly aligns with the missing geometric state. The \textbf{Sterile Vacancy} occupies the unanchored Right-Handed slot. Therefore, the framework predicts that all experimental searches for heavy, propagating right-handed neutrinos (such as those proposed by the Seesaw Mechanism) will yield null results. They are not physical particles; they are the required static reference capacity of the network.

\subsection{The 16-Channel Hardware Limit: Exhausting the Substrate}
\label{sec:enumeration_table}

By combining the active protocols with the sterile vacancy, we can exhaustively map the Standard Model fermion architecture to the $\nu=16$ hardware channels. The channels follow a strict \textbf{Hierarchy of Geometric Decoupling}, ordered from the most densely anchored protocols down to the completely uncoupled reference state:

\begin{enumerate}
    \item \textbf{Color + Active Boundary (Channels 1--6):} The heaviest geometric footprint. These protocols route through the $N_{res}=3$ internal color bulk and couple to the active Left-Handed temporal boundary (forming the $SU(2)_L$ Quark Doublets).
    \item \textbf{Color + Static Boundary (Channels 7--12):} These protocols route through the color bulk but are decoupled from the active temporal boundary, residing in the Right-Handed internal sector (Quark Singlets).
    \item \textbf{No Color + Active Boundary (Channels 13--14):} These protocols bypass the color bulk entirely but remain coupled to the active Left-Handed temporal boundary (forming the Lepton Doublet: $\nu_{eL}$ and $e_L$).
    \item \textbf{No Color + Static Boundary (Channel 15):} Bypasses the color bulk and decoupled from the active temporal boundary (the Right-Handed Electron, $e_R$).
    \item \textbf{Total Decoupling (Channel 16):} Bypasses the color bulk, decoupled from the active boundary, and possesses zero topological anchors (The Sterile Vacancy).
\end{enumerate}

\begin{table}[h]
\centering
\caption{The $\nu=16$ chiral channels of the $E_8 \to D_4 \oplus D_4$ substrate. The 15 active protocols perfectly exhaust the routing possibilities. Generational copies (Tier 1, Tier 2, Tier 3) are simply higher-dimensional temporal overtones of these exact same 16 hardware channels.}
\label{tab:channels}
\begin{tabular}{|c|c|c|c|c|c|c|c|c|}
\hline
\textbf{Ch.} & \textbf{Protocol} & \textbf{Color Bulk} & \textbf{Isospin ($I_3$)} & \textbf{Hand} & \textbf{Charge ($Q$)} & \textbf{Anchor} & \textbf{Tier 0} \\
\hline
1--3 & Up-Type & $N_{\rm res}=3$ & $+\frac{1}{2}$ & L & $+\frac{2}{3}$ & $D=4$ & $u_L$ \\
4--6 & Down-Type & $N_{\rm res}=3$ & $-\frac{1}{2}$ & L & $-\frac{1}{3}$ & $D=4$ & $d_L$ \\
7--9 & Up-Type & $N_{\rm res}=3$ & 0 & R & $+\frac{2}{3}$ & $D=4$ & $u_R$ \\
10--12 & Down-Type & $N_{\rm res}=3$ & 0 & R & $-\frac{1}{3}$ & $D=4$ & $d_R$ \\
13 & Neutrino & None & $+\frac{1}{2}$ & L & 0 & None & $\nu_{eL}$ \\
14 & Electron & None & $-\frac{1}{2}$ & L & $-1$ & $\chi=2$ & $e_L$ \\
15 & Electron & None & 0 & R & $-1$ & $\chi=2$ & $e_R$ \\
16 & Sterile & None & 0 & --- & 0 & None & $\nu_R$ \\
\hline
\end{tabular}
\end{table}

\subsection{Protocol Resonance Tiers}

Each of the four protocols can execute the same base protocol at \textbf{distinct resonant tiers} dictated by their geometric dimensionality ($\Delta^1, \Delta^2, \Delta^3$). In standard physics, this is known as the three “generations” seen in the Lepton and Quark elementary particles.

The substrate has constraints that must be taken into account to successfully persist at different tiers.  The specific thermodynamic corrections arising from capacity limits, boundary leakage, and isospin structure. Thus the \textbf{mass} of an elementary particle is the thermodynamic latency of maintaining its specific dimensional standing wave:

\paragraph{The Prohibition of a Fourth Generation:}
Why does the Standard Model stop at three generations? In standard physics, this is an empirical mystery. In the $E_8$-Persistence framework, it is a strict capacity limit. The generations correspond to spatial dimensions: 0D, 1D, 2D, and 3D. The Top Quark ($\Delta^3$) hits the absolute volumetric capacity ceiling of the $D=4$ local node. A hypothetical fourth generation would require a 4D standing wave ($\Delta^4$), which fundamentally exceeds the geometric boundaries of the node. Such a state would trigger instantaneous Causal Aliasing and dissolve. Thus, exactly three heavy generations can exist.

\paragraph{Dissolving the Hierarchy Problem:}
Standard physics views the massive scale difference between the Electron and the Top Quark ($\sim 340,000\times$) as an unnatural fine-tuning problem. Geometrically, this hierarchy is trivial. The Electron is a 0D point ($\Delta^0=1$). The Top Quark is a 3D volume ($\Delta^3 \approx 79,507$) subjected to a 4D capacity ceiling multiplier ($D+1/D = 4.25$). The mass gap is simply the geometric volume of the lattice resonance ($79507 \times 4.25 \approx 337,900$).
